%% =========================================================================
%% Cover Page and Abstract
%% =========================================================================

\begin{titlepage}
    \centering
    \vspace*{1.0cm}
    
    {\Huge \bfseries \color{navyblue} Early-Life Disadvantage and Later-Life Mental and Cognitive Health in India\par}
    
    \vspace{0.6cm}
    {\LARGE \bfseries \color{darkslate} Evidence from the Longitudinal Ageing Study in India (LASI Wave 1)\par}
    
    \vspace{0.4cm}
    {\large \itshape A Survey-Weighted Epidemiological Investigation with Continuous Cognitive Sensitivity Analysis and Health Policy Implications\par}
    
    \vspace{1.5cm}
    \rule{0.85\textwidth}{1.5pt}
    \vspace{1.5cm}
    
    \begin{minipage}[t]{0.85\textwidth}
        \large
        \textbf{Course / Programme:} SOC450: Sociology of Ageing and Health / Social Epidemiology \\
        \textbf{Academic Term:} Autumn Semester (Semester V), 2026--2027 \\
        \textbf{Department:} Department of Humanities and Social Sciences \\
        \textbf{Institution:} Indian Institute of Technology / University \\
        \textbf{Faculty Supervisor:} Prof. [Faculty Supervisor Name]
    \end{minipage}
    
    \vspace{1.8cm}
    
    \textbf{\large Submitted by: Project Group Members}
    
    \vspace{0.4cm}
    
    \begin{table}[H]
        \centering
        \large
        \begin{tabular}{llr}
            \toprule
            \textbf{S. No.} & \textbf{Student Name} & \textbf{Roll Number} \\
            \midrule
            1 & Harshit & [Roll Number 1] \\
            2 & [Group Member 2] & [Roll Number 2] \\
            3 & [Group Member 3] & [Roll Number 3] \\
            4 & [Group Member 4] & [Roll Number 4] \\
            \bottomrule
        \end{tabular}
    \end{table}
    
    \vfill
    
    {\large \textbf{Date of Submission:} October 2026\par}
    \vspace*{0.5cm}
\end{titlepage}

\newpage
\pagenumbering{roman}
\setcounter{page}{2}

\begin{center}
    {\Large \bfseries \color{navyblue} Abstract}
\end{center}

\vspace{0.5cm}

\noindent \textbf{Problem Statement:} 
Whether childhood socioeconomic hardships and infectious illnesses leave permanent scars on cognitive vitality and psychological well-being in later life remains a central question in life-course epidemiology. In low- and middle-income countries like India, investigating these relationships is complicated by widespread illiteracy, selective survival, and potential testing artifacts in cognitive assessments.

\vspace{0.3cm}
\noindent \textbf{Objective:} 
To evaluate the associations between recalled childhood socioeconomic disadvantage, general health, and bedridden illness with later-life depressive symptoms and cognitive performance among older Indian adults, and to examine whether these associations persist across adult socioeconomic attainment and stages of ageing.

\vspace{0.3cm}
\noindent \textbf{Dataset:} 
Microdata from Wave 1 of the Longitudinal Ageing Study in India (LASI, 2017--2018), analyzing a nationally representative sample of $66,606$ community-dwelling adults aged 45 years and older across all 36 States and Union Territories.

\vspace{0.3cm}
\noindent \textbf{Methodology:} 
We estimated survey-weighted modified Poisson regressions with design-based stratified-cluster sandwich variance estimators (72 strata, 2,428 clusters, $df=2,356$) on common complete-case samples ($N=62,047$ for depression; $N=57,245$ for cognition). Depression was defined using the 10-item CES-D ($\ge 4.0$). To account for differential test administration in populations with high illiteracy, the primary cognitive outcome was defined using an education-fair memory and orientation composite (0--28 points; weighted lowest decile $\le 11.0$). Age-graded effect modification was formally assessed via joint Wald $F$ tests, and domain-specific continuous scores were examined in an extension.

\vspace{0.3cm}
\noindent \textbf{Principal Findings:} 
Depressive symptoms (prevalence 28.2\%) demonstrated persistent associations with early disadvantage that resisted attenuation by adult schooling or wealth. Adults reporting poor childhood health had a 20\% higher prevalence of depression (Adjusted Prevalence Ratio $[\text{APR}] = 1.20$, 95\% CI: 1.07--1.34), poor childhood SES had a 21\% higher prevalence ($\text{APR} = 1.21$, 95\% CI: 1.03--1.42), and childhood bedridden illness had a 16\% higher prevalence ($\text{APR} = 1.16$, 95\% CI: 1.05--1.28). In contrast, cognitive impairment disparities attenuated fully after adjusting for adult education and consumption expenditure ($\text{APR} = 0.93$, 95\% CI: 0.82--1.05). Age interactions were statistically null for depression ($F=0.34, p=0.851$), indicating uniform lifelong vulnerabilities rather than age-widening divergence. Continuous extension models confirmed that residual cognitive deficits were restricted strictly to memory recall rather than orientation.

\vspace{0.3cm}
\noindent \textbf{Conclusion:} 
Early-life health and economic precarity leave enduring psychiatric footprints across the adult life span in India. Conversely, late-life cognitive differentials reflect educational access and test construction. Public health policies must integrate early childhood health investments with universal education and community-based geriatric mental health screening.

\vspace{0.5cm}
\noindent \textbf{Keywords:} Life-Course Epidemiology, Cognitive Ageing, Late-Life Depression, LASI Wave 1, Survey-Weighted Poisson Regression, Social Determinants of Health, India.

\newpage
\tableofcontents
\newpage
\pagenumbering{arabic}
\setcounter{page}{1}
